\documentclass[letterpaper]{article}
\usepackage[margin=2cm]{geometry}
\usepackage[english]{babel}
\usepackage{amsfonts}
\usepackage{amsthm}
\usepackage{amsmath}
\usepackage{amssymb}
\usepackage[edges]{forest}
\usepackage{tikz}
\pagestyle{empty}
\usepackage{caption}
\usepackage{enumitem}
\usepackage{setspace}
\usepackage{fdsymbol}
\usepackage{pgfplots}
\usepackage{cancel}
\usepackage{soul}
\usepackage{xcolor}
\usepackage{listings}
\usepackage{graphicx}

\pgfplotsset{compat=1.18} 

\begin{document}
\flushleft
\renewcommand{\labelitemiii}{\(\smallblackcircle\)}
\renewcommand{\labelitemii}{\(\smallblackcircle\)}

\lstset{
  basicstyle=\ttfamily,
  mathescape
}

\section*{A2 MODELS}

\subsection*{Q1.}

Define the TM M as follows:

\begin{align*}
   & \Sigma = \{0, 1\}                          \\
   & S = \{ s_{start}, s_0, s_1, s_{end}\}      \\
   & \delta (s_{start}, 0)= (s_0, 0, 1)         \\
   & \delta (s_{start}, 1) = (s_1, 0, 1)        \\
   & \delta (s_{start}, \beta)= (s_{end}, 0, 1) \\
   & \delta (s_{0}, 0)= (s_0, 0, 1)             \\
   & \delta (s_{1}, 0)= (s_0, 1, 1)             \\
   & \delta (s_{0}, 1) = (s_1, 0, 1)            \\
   & \delta (s_{1}, 1) = (s_1, 1, 1)            \\
   & \delta (s_{0}, \beta)= (s_{end}, 0, 1)     \\
   & \delta (s_{1}, \beta)= (s_{end}, 1, 1)     \\
\end{align*}

\subsubsection*{Complexity}
The transition function always moves right and stops once it sees $\beta$.
So the runtime is $n+1=\mathcal{O} (n)$ for any $x \in \{0, 1\}^n$.

\subsubsection*{Correctness}
\begin{itemize}
  \item First, M writes 0 in the first cell. If the original value was 1 or 0,
        then M saves it to the state and moves right. If the original value was
        $\beta$ then M simply halts.
  \item By induction on $i=1,...,n-1$, after $i$ steps, first $i$ cells of
        the worktape contain $0 x_1 x_2 ... x_{i-1}$ and the state is $s_{x_i}$.
        In step $i+1$, M writes $x_i$ in the $x+1$ cell and saves $x+1$ in the state.
  \item In step $n$, the machine state is $s_{x_{n-1}}$ so M writes
        $x_{n-1}$, reads $\beta$ and halts.
\end{itemize}

\subsection*{Q2.}

Define TM as follows:
\begin{itemize}
  \item Mark the first bit as 1'. (Assuming that there are no leading 0's)
  \item Go all the way right
  \item Go 1 step left
        \begin{itemize}
          \item If it sees 1, set to 0 and go 1 step left. Repeat this until it sees a 0.
          \item If we see a 0, set it to 1 and halt.
        \end{itemize}
  \item If we reach 1' without seeing a 0
        \begin{itemize}
          \item Leave first bit as 1
          \item Walk right until we see a $\beta$, write 0 and halt.
        \end{itemize}
\end{itemize}

\subsection*{Q3.}

Define TM as follows:
\begin{itemize}
  \item For each unmarked bit $\ell$
        \begin{itemize}
          \item Mark is as seen ($\ell'$) and set state to just seen $\ell$ ($s_\ell$)
          \item Walk right to the rightmost unmarked symbol $r$ and mark is as seen ($r'$)
                \begin{itemize}
                  \item If there is no such symbol, then accept and halt.
                  \item Check that $r = \ell$
                        \begin{itemize}
                          \item If $r \ne \ell$, reject and halt.
                          \item If equal, continue
                        \end{itemize}
                \end{itemize}
        \end{itemize}
  \item Once all bits are marked, accept and halt.
\end{itemize}

\subsubsection*{Complexity}

M processes at most $\lceil n/2 \rceil$ bits on the left. For each of these
bits, M walks $\le n$ bits right and back. $\Rightarrow \mathcal{O}(n^2)$.

\subsubsection*{Q4.}

Let M be a TM with special end states $s_{accept}$ and $s_{reject}$.

Define M' as follows:
\begin{itemize}
  \item Runs the same steps as M
  \item Except, when M accepts or rejects, M' instead:
        \begin{itemize}
          \item Goes all the way right
          \item Erases the tape, right to left
          \item In the leftmost cell, writes 1 for accept and 0 for reject.
        \end{itemize}
\end{itemize}

\subsubsection*{Complexity}

For all intermediate steps, M' runs in the same time as M $\Rightarrow T(n)$.
In the last step, M' walks up to the entire length of the tape twice.
There are at most $T(n)$ cells to traverse and erase. Therefore, M' runs in
time $\mathcal{O}(T(n))$

\subsection*{Q5.}

Let M be a two-tape TM running in time $T(n)$.

Let M' be a standard TM containing $w_1 \# w_2$ on the tape.
Mark the first bit of $w_1$ and $w_2$ as the head with '.
Simulate each action of M as follows:
\begin{itemize}
  \item Walk leftmost then right to the head on $w_1$.
        Simulate action in $w_1$, unmarking the previous head and marking
        the new head.
  \item Walk right (passed the delimiter) to the head on $w_2$.
        Simulate action on $w_2$, unmarking the previous head and marking
        the new head.
  \item If $w_1$ runs out of space, shift entire contents of $\# w_2$
        to the right and pad with $\beta$.
\end{itemize}

\subsubsection*{Complexity}

For each action by M, M' walks up to $T(n)$ cells to get to the next correct
head location, simulates the action, and possibly shifts the tape contents
1 step to the right which is at most $T(n)$ in length.

$\Rightarrow$ M' simulates M in time $\mathcal{O}(T(n)^2) = poly(T(n))$.

\subsubsection*{Correctness}

TODO

\subsection*{Q6.}

Let M be a k-tape TM running in $T(n)$ time.
Let M' be a standard TM, with the tape content formatted as
$w_1 \# w_2 \# ... \# w_k $. i.e. each $w_i$ separated by $\#$.
Mark the first bit of each $w_i$, indicating the head location.

Simulate M with M' as follows:
\begin{itemize}
  \item Walk leftmost then right to the head on $w_1$. Simulate action on $w_1$
        and update the head location by marking the new head.
  \item Walk right until the nearest delimiter, then walk to the head on $w_2$.
        Simulate the action on $w_2$, and update the head location.
  \item Walk right until the nearest delimiter, then walk to the head on $w_3$.
        Simulate the action on $w_3$, and update the head location.
  \item ...
  \item Walk right until the last delimiter, then walk to the head on $w_k$.
        Simulate the action on $w_k$, and update the head location.
  \item Repeat this process. If any of the $w_i$ run out of space, shift the
        entire right side contents to the right.
\end{itemize}

\subsubsection*{Complexity}

For each action by M, M walks at most $k \cdot T(n)$ steps to get to the
head of the next section to work on, simulates the action, and possibly
shifts at most $T(n)$ cells right 1 step. $\Rightarrow $ M' runs in
$\mathcal{O}(T(n)^2) = poly(T(n))$

\subsubsection*{Correctness}

TODO

\subsection*{Q7.}

Let M be a two-way-infinite tape.
Let M' be a one-way-infinite tape.
Simulate M with M' as follows:
\begin{itemize}
  \item Mark the leftmost cell with a delimiter.
  \item Copy the input on the right of the delimiter.
  \item For each action of M
        \begin{itemize}
          \item If M moves right, M' also moves right.
          \item If M moves left, M' also moves left.
          \item If M' every reads the delimiter, then it shifts the entire
                tape's content from $\# x_1 x_2 ...$ to $\# \beta x_1 x_2 ...$.
                Then it walks right until it reaches the delimiter,
                moves one step right and continues
                the simulation.
        \end{itemize}
\end{itemize}

\subsubsection*{Complexity}

For each action of M, M' does the same thing, except, if M reads the delimiter,
it shift the entire tape contents one cell to the right. So M' walks at most
$T(n)$ cells to the right, and then up to $T(n)$ cells to the left.

$\Rightarrow$ M' runs in time $\mathcal{O}(T(n)^2) = poly(T(n))$.

\subsubsection*{Correctness}

\subsection*{Q8.}

Let $k$ be the size of $\Sigma$. Then, each symbol is encoded by
$\lceil \log \lvert \Sigma \rvert \rceil$ bits.

Define TM M' as follows:
\begin{itemize}
  \item For each sequence of bits representing a symbol in $\Sigma$, define
        states for each prefix. i.e. if the encoding of a symbol is 1010, then
        create states $s_{1}, s_{10}, s_{101}, s_{1010}$ for just saw this sequence.
        Once the entire symbol sequence is recognized, it can write the corresponding
        bits of the output symbol from right to left (determined by the states).
        Then, if M moves left, M' moves $\lceil \log \lvert \Sigma \rvert \rceil$
        steps left. Or if M moves right, M' moves $\lceil \log \lvert \Sigma \rvert \rceil$
        steps right.
  \item For each action of M, M' performs
        $3 \cdot \lceil \log \lvert \Sigma \rvert \rceil$
        times as many steps, so it computes $f$ in time
        $\mathcal{O}(3 \cdot \lceil \log \lvert \Sigma \rvert \rceil \cdot T(n)) = \mathcal{O}(T(n))$
\end{itemize}

\subsubsection*{Complexity}

TODO

\subsubsection*{Correctness}

TODO

\subsection*{Q9.}

\setlength{\parskip}{10pt}{
  Consider $L = \{x : x \text{ is a palindrome}\}$.

  To decide if a string is a palindrome, the TM must look at all
  $n$ bits if $n$ is even or $n-1$ bits if $n$ is odd
  (middle bit can be ignored).

  For the sake of contradiction, suppose that $L$ can be decided in
  $o(n)$. Then, for any input, there must be at least one bit
  (that is not the center bit when $n$ is odd)
  that the TM does not look at.

  Consider a palindrome string $x_1$.
  There is at least one bit in $x_1$ that the TM does not look at
  (and is not the middle bit if the string has odd length).
  Let this bit be in position $i$.
  Let $x_2$ be $x_1$ with the $i$th bit flipped. Thus $x_2$ is not a palindrome.
  The TM accepts both $x_1$ and $x_2$ (since it does not look at the $i$th bit).
  This is a contradiction.

  Therefore, $L$ cannot be decided in $o(n)$.
}

\subsection*{Q10.}

Consider $f(x) = 0^{2^{2^{2^{2^n}}}}$. i.e. $f$ takes any binary string of
length $n$ as input and outputs 0 $2^{2^{2^{2^n}}}$ times. This function
cannot be computed in time $o(2^{2^{2^{2^n}}})$ since it takes at least 
$2^{2^{2^{2^n}}}$ writes to write $2^{2^{2^{2^n}}}$ 0's. 

\end{document}
